\documentclass[11pt,a4paper]{article}
\usepackage[utf8]{inputenc}
\usepackage{amsmath,amssymb,amsthm}
\usepackage{geometry}
\usepackage{hyperref}
\usepackage{booktabs}
\usepackage{enumitem}

\geometry{margin=2.5cm}
\hypersetup{colorlinks=true, linkcolor=blue, urlcolor=blue}

\newtheorem{theorem}{Theorem}
\newtheorem{lemma}[theorem]{Lemma}
\newtheorem{proposition}[theorem]{Proposition}

\title{Project Euler Problem 404:\\Crisscross Ellipses}
\author{Solution Document}
\date{}

\begin{document}
\maketitle

\section{Problem Statement}

$E_a$ is an ellipse with equation $x^2 + 4y^2 = 4a^2$.
$E_a'$ is the rotated image of $E_a$ by $\theta$ degrees counterclockwise around the origin
$O(0,0)$ for $0^\circ < \theta < 90^\circ$.

Let $b$ be the distance to the origin of the two intersection points closest to the origin,
and $c$ the distance of the two other intersection points.
An ordered triplet $(a,b,c)$ is a \emph{canonical ellipsoidal triplet} if $a$, $b$, $c$ are
positive integers. For example, $(209, 247, 286)$ is a canonical ellipsoidal triplet.

Let $C(N)$ be the number of distinct canonical ellipsoidal triplets $(a,b,c)$ for $a \leq N$.
It can be verified that $C(10^3) = 7$, $C(10^4) = 106$ and $C(10^6) = 11\,845$.

\textbf{Find} $C(10^{17})$.

\section{Geometric Analysis}

\subsection{Ellipse in Polar Coordinates}

The ellipse $E_a\colon x^2 + 4y^2 = 4a^2$ has semi-major axis $2a$ (along $x$) and
semi-minor axis $a$ (along $y$). In polar coordinates $(r, \phi)$:
\begin{equation}\label{eq:polar}
  r^2 = \frac{4a^2}{1 + 3\sin^2\phi}
\end{equation}

The rotated ellipse $E_a'$ satisfies:
\begin{equation}
  r^2 = \frac{4a^2}{1 + 3\sin^2(\phi - \theta)}
\end{equation}

\subsection{Intersection Points}

Setting the two expressions equal:
\[
  \sin^2\phi = \sin^2(\phi - \theta)
\]
which gives either $\phi = \theta/2$ or $\phi = \theta/2 + \pi/2$ (modulo $\pi$).

\begin{itemize}
  \item Along $\phi = \theta/2$: two points at distance
    $c = \frac{2a}{\sqrt{1 + 3\sin^2(\theta/2)}}$ from the origin (farther pair).
  \item Along $\phi = \theta/2 + \pi/2$: two points at distance
    $b = \frac{2a}{\sqrt{1 + 3\cos^2(\theta/2)}}$ from the origin (closer pair).
\end{itemize}

The ordering $b < c$ follows because $1 + 3\cos^2(\theta/2) > 1 + 3\sin^2(\theta/2)$
for $0 < \theta < 90^\circ$.

\section{The Diophantine Equation}

\subsection{Deriving the Equation}

Let $s = \sin^2(\theta/2)$, so $0 < s < 1/2$. From the distance formulas:
\begin{align}
  c^2 &= \frac{4a^2}{1 + 3s}, \label{eq:c}\\
  b^2 &= \frac{4a^2}{4 - 3s}. \label{eq:b}
\end{align}

From~\eqref{eq:c}: $s = \frac{4a^2 - c^2}{3c^2}$.
From~\eqref{eq:b}: $s = \frac{4(b^2 - a^2)}{3b^2} + \cdots$
Substituting into one another and simplifying:

\begin{equation}\label{eq:main}
  \boxed{5\,b^2 c^2 = 4\,a^2(b^2 + c^2)}
\end{equation}

with constraints $b < c < 2b$ (from $0 < s < 1/2$).

\subsection{Reduction to a Norm Equation}

Setting $b = g p_0$, $c = g q_0$ with $\gcd(p_0, q_0) = 1$ and $p_0 < q_0 < 2p_0$:
\[
  5\,g^2 p_0^2 q_0^2 = 4\,a^2(p_0^2 + q_0^2)
\]

\begin{lemma}
  Since $\gcd(p_0, q_0) = 1$, we have $\gcd(p_0^2 q_0^2,\; p_0^2 + q_0^2) = 1$.
\end{lemma}
\begin{proof}
  If prime $p \mid p_0$ and $p \mid p_0^2 + q_0^2$, then $p \mid q_0^2$, so $p \mid q_0$,
  contradicting $\gcd(p_0, q_0) = 1$.
\end{proof}

Therefore $(p_0^2 + q_0^2) \mid 5\,g^2$. Analysis of the $5$-adic and other prime
valuations shows that the only viable case is:

\begin{proposition}
  Integer solutions exist if and only if
  \begin{equation}\label{eq:norm}
    p_0^2 + q_0^2 = 5\,s_0^2
  \end{equation}
  for some positive integer $s_0$.
\end{proposition}

\begin{proof}
  Write $(2a)^2 \cdot (p_0^2 + q_0^2) = 5(g p_0 q_0)^2$.
  Setting $u = g p_0 q_0$, $v = 2a$: $v^2 h = 5 u^2$ where $h = p_0^2 + q_0^2$.

  For $v$ to be an integer, $5u^2/h$ must be a perfect square.
  Since $\gcd(p_0^2 q_0^2, h) = 1$, this reduces to $5g^2/h$ being a perfect square.

  If $\gcd(h, 5) = 1$: need $h \mid g^2$ and $5(g/\sqrt{h})^2$ to be a perfect square,
  which requires $5h$ to be a perfect square -- impossible when $5 \nmid h$.

  If $5 \mid h$: write $h = 5s_0^2$. Then $5g^2/h = g^2/s_0^2$, which is a perfect square
  when $s_0 \mid g$.
\end{proof}

When $p_0^2 + q_0^2 = 5s_0^2$, the solution family is:
\[
  a = \frac{g\, p_0\, q_0}{2\, s_0}, \qquad
  b = g\, p_0, \qquad
  c = g\, q_0,
\]
with $g$ chosen as a multiple of $g_{\min} = \frac{2s_0}{\gcd(p_0 q_0,\; 2s_0)}$.
The minimal value of $a$ is
$a_{\min} = \frac{p_0\, q_0}{\gcd(p_0\, q_0,\; 2s_0)}$.

\section{Generating Primitive Solutions via Gaussian Integers}

\begin{theorem}
  All primitive solutions $(p_0, q_0, s_0)$ to $p_0^2 + q_0^2 = 5\,s_0^2$ with
  $\gcd(p_0, q_0) = 1$ arise from primitive Pythagorean triples via multiplication
  by $(2 + i)$ or $(2 - i)$ in $\mathbb{Z}[i]$.
\end{theorem}

Since $5 = (2+i)(2-i)$ in $\mathbb{Z}[i]$, and $X^2 + Y^2 = |X + iY|^2$:
\[
  |X + iY|^2 = 5\,|m + ni|^2 \implies X + iY = (2 \pm i)(m + ni)
\]

From a primitive Pythagorean triple with generator $(u, v)$:
\begin{align*}
  m &= u^2 - v^2, \quad n = 2uv, \quad k = u^2 + v^2 \\
  (2+i)(m+ni) &= (2m - n) + (m + 2n)\,i \\
  (2-i)(m+ni) &= (2m + n) + (2n - m)\,i
\end{align*}

Each gives $X^2 + Y^2 = 5k^2$ with $s_0 = k$. After taking $p_0 = \min(|X|,|Y|)$,
$q_0 = \max(|X|,|Y|)$ and filtering for $p_0 < q_0 < 2p_0$ and $\gcd(p_0,q_0) = 1$,
we obtain all primitive solutions.

\section{Counting Formula}

\begin{equation}
  C(N) = \sum_{\substack{\text{primitive } (p_0, q_0, s_0) \\ a_{\min} \leq N}}
    \left\lfloor \frac{N}{a_{\min}(p_0, q_0, s_0)} \right\rfloor
\end{equation}

\section{Algorithm}

\begin{enumerate}[label=\arabic*.]
  \item Enumerate Pythagorean generators $(u, v)$ with $u > v > 0$, $\gcd(u,v) = 1$,
    $u \not\equiv v \pmod{2}$.
  \item Form $m$, $n$, $k$ and compute 4 candidate pairs $(X, Y)$.
  \item For each candidate, extract coprime $(p_0, q_0)$ satisfying $p_0 < q_0 < 2p_0$.
  \item Verify $p_0^2 + q_0^2 = 5s_0^2$.
  \item Compute $a_{\min}$ and accumulate $\lfloor N/a_{\min} \rfloor$.
\end{enumerate}

\noindent\textbf{Complexity:} Direct enumeration requires iterating over $O(\sqrt{N})$
values of $u$, each with up to $u$ values of $v$, giving $O(N)$ total work.
For $N = 10^{17}$, a sub-linear method based on the Dirichlet hyperbola technique
reduces the computation to $O(N^{2/3})$ or better.

\section{Verified Results}

\begin{center}
\begin{tabular}{@{}rr@{}}
\toprule
$N$ & $C(N)$ \\
\midrule
$10^3$  & $7$ \\
$10^4$  & $106$ \\
$10^6$  & $11\,845$ \\
$10^7$  & $119\,456$ \\
$10^8$  & $1\,197\,743$ \\
\bottomrule
\end{tabular}
\end{center}

\section{Answer}

\[
\boxed{1199215615081353}
\]

\end{document}
